\section{Quantum confinement in the homogeneous gauge--scalar model}
\label{qh:chapter}

The object quantized here is a precisely specified interacting
matrix mechanics: three anti-Hermitian traceless gauge matrices,
one Hermitian traceless adjoint scalar, and one complex triplet.
It is obtained by spatially constant substitution in the gauge--scalar
action used in the quaternionic continuation. All of these fields
remain variables. In particular the triplet is not constrained to a
four-real-dimensional slice or to a scalar minimum curve.

The conclusions are a compact-resolvent Hamiltonian, a unique
spatially constant $SU(3)$-invariant ground state, and a strictly
positive first excitation gap at fixed positive volume and couplings.
The explicit form inequality \eqref{qh:full-form-bound} is the
quantitative step. It combines the exact $SU(3)$ root frequencies
with the full quartic scalar potential. Bosonic commutator-matrix
discreteness is a classical result; the primary attribution and
the complete proof in the retained metric are given below.
The present use is to answer a concrete quantization question for
the already specified adjoint-plus-triplet action, not to claim
a new general discreteness mechanism.

For the classical comparison the spacetime signature is $(+---)$
and the original action density is
\[
 \mathcal L=\frac1{2g^2}\operatorname{tr}F_{\mu\nu}F^{\mu\nu}
       +\frac12\operatorname{tr}(D_\mu\Phi D^\mu\Phi)
       +(D_\mu b)^*D^\mu b-V(\Phi,b),
\]
where
\[
 F_{\mu\nu}=\partial_\mu A_\nu-\partial_\nu A_\mu+[A_\mu,A_\nu],
 \quad D_\mu\Phi=\partial_\mu\Phi+[A_\mu,\Phi],
 \quad D_\mu b=\partial_\mu b+A_\mu b.
\]
The potential is written in full below. The finite-model theorem
does not require the construction of a compact complex structure
on $S^6$, and it neither verifies nor refutes that separate global
construction. The precise classical map to the torus field theory,
its noninjective map on gauge orbits, and an explicit obstruction
to one proposed quantum-subsector map are proved at the end.

% Insertion-ready proof; no canonical source is edited by this artifact.
% Required packages/environments: amsmath, amssymb, standard theorem/proof.
\subsection{The retained configuration space and action}

Fix a flat spatial torus of volume \(\mathcal V_3>0\), a gauge coupling
\(g^2>0\), and \(\hbar>0\). Retain \(m^2\geq0\), \(\Gamma>0\) and
\(\kappa\in\mathbb R\). We quantize the finite-dimensional homogeneous
mechanics, not a claimed invariant quantum sector of the full field theory.
The configuration space, with no quotient already performed, is
\[
 \mathcal Q=\mathfrak{su}(3)^3\oplus
       \operatorname{Herm}_3(\mathbb C)_0\oplus\mathbb C^3,
 \qquad \dim_{\mathbb R}\mathcal Q=24+8+6=38.
\]
Its elements are \(q=(A_1,A_2,A_3,\Phi,b)\). On \(\mathfrak{su}(3)\)
use \((X,Y)_{\mathfrak g}=-\operatorname{tr}(XY)\); the scalar metric is
exactly \(\operatorname{tr}(d\Phi^2)+2\,db^*db\).
The trace-free condition on \(\Phi\) is part of the linear state space.

For a homogeneous temporal connection \(A_0\), set
\[
 E_i=\dot A_i+[A_0,A_i],\quad
 X=\dot\Phi+[A_0,\Phi],\quad v=\dot b+A_0b,\quad
 s=\operatorname{tr}\Phi^2+2b^*b.
\]
The spatial integral of the retained \((+---)\), anti-Hermitian density is
\begin{align}
 L_{\rm hom}&=\mathcal V_3\left[
   \frac1{g^2}\sum_i\|E_i\|_{\mathfrak g}^2
      +\frac12\operatorname{tr}X^2+v^*v-U(q)\right],
       \label{qh:hq-L}\\
 U(q)&=\frac1{g^2}\sum_{i<j}\|[A_i,A_j]\|_{\mathfrak g}^2
       +\frac12\sum_i\operatorname{tr}[A_i,\Phi]^2
       +\sum_i\|A_i b\|^2+V(\Phi,b),\notag\\
 V(\Phi,b)&=\frac{m^2}{2}s+\frac{\Gamma}{4}s^2
                      -\kappa(\det\Phi-b^*\Phi b).\notag
\end{align}
Here \([A_i,\Phi]\) is Hermitian, so its displayed squared trace is
nonnegative. There is no velocity for \(A_0\).

\begin{proposition}[Canonical momenta, Hamiltonian and exact Gauss map]
\label{qh:hq-canonical}
Use the canonical one-form
\[
 \vartheta=\sum_i-\operatorname{tr}(\Pi_i\,dA_i)
       +\operatorname{tr}(\Pi_\Phi\,d\Phi)
       +2\operatorname{Re}(\pi_b^*db).
\]
The momenta and Hamiltonian are
\begin{align}
 \Pi_i&=\frac{2\mathcal V_3}{g^2}E_i,\qquad
 \Pi_\Phi=\mathcal V_3X,\qquad \pi_b=\mathcal V_3v,\notag\\
 H_{A_0}&=H_0-\mathcal M_{A_0},\label{qh:hq-H}\\
 H_0&=\frac{g^2}{4\mathcal V_3}\sum_i\|\Pi_i\|_{\mathfrak g}^2
       +\frac1{2\mathcal V_3}\operatorname{tr}\Pi_\Phi^2
       +\frac1{\mathcal V_3}\pi_b^*\pi_b+\mathcal V_3U(q).\notag
\end{align}
For \(\xi\in\mathfrak{su}(3)\), write
\[
 K_\xi(q)=([\xi,A_i]_{i=1}^3,[\xi,\Phi],\xi b),\qquad
 \mathcal M_\xi=\vartheta(K_\xi)=(\mathcal G,\xi)_{\mathfrak g}.
\]
The exact matrix-valued momentum map is
\begin{equation}\label{qh:hq-Gauss}
 \mathcal G=\sum_i[A_i,\Pi_i]-[\Phi,\Pi_\Phi]
          -(b\pi_b^*-\pi_b b^*)_{\rm tf}
 =\mathcal V_3\left(\frac2{g^2}\sum_i[A_i,E_i]-(J_0)_{\rm tf}\right),
\end{equation}
where \(J_0=[\Phi,X]+bv^*-vb^*\).
Thus \(\mathcal G=0\) is precisely the full homogeneous Gauss equation.
\end{proposition}
\begin{proof}
Differentiate \eqref{qh:hq-L} using the stated duality pairings.
In particular, differentiation of \(v^*v\) gives
\(2\mathcal V_3\operatorname{Re}(v^*\delta\dot b)\), explaining
\(\pi_b=\mathcal V_3v\) rather than \(2\mathcal V_3v\) in this complex
duality convention. The velocity is the covariant velocity minus
\(K_{A_0}\); its Legendre transform is consequently \(H_0-\vartheta(K_{A_0})\).
The ordinary quadratic Legendre transform gives every coefficient of
\eqref{qh:hq-H}. Cyclicity gives
\[
 -\operatorname{tr}(\Pi_i[\xi,A_i])
       =-\operatorname{tr}([A_i,\Pi_i]\xi),\quad
 \operatorname{tr}(\Pi_\Phi[\xi,\Phi])
       =\operatorname{tr}([\Phi,\Pi_\Phi]\xi),
\]
while anti-Hermiticity gives
\[
 2\operatorname{Re}(\pi_b^*\xi b)
       =\operatorname{tr}((b\pi_b^*-\pi_b b^*)\xi).
\]
These identities prove the first expression for \(\mathcal G\);
substitution of the momenta proves the second. Variation of the
independent \(A_0\) multiplier imposes \(\mathcal G=0\), even after a
temporal gauge is selected for the evolution. The primary momentum
conjugate to \(A_0\) is zero; \(A_0\) is not an additional coordinate
of \(\mathcal Q\). Since \(\mathcal M_\xi=p\cdot K_\xi\), canonical
Poisson brackets give
\(\{\mathcal M_\xi,\mathcal M_\eta\}=\mathcal M_{[\xi,\eta]}\).
Gauge invariance of every trace and norm in \(H_0\) also gives
\(\{H_0,\mathcal M_\xi\}=0\).
\end{proof}

\paragraph{Coordinates and measure without changing the metrics.}
Choose any fixed orthonormal basis \(T_a\), \(1\leq a\leq8\), for
\(-\operatorname{tr}\) on \(\mathfrak{su}(3)\), and put \(S_a=iT_a\).
Then \(\operatorname{tr}(S_aS_b)=\delta_{ab}\). Write
\[
 A_i=\sum_a a_{ia}T_a,\quad
 \Phi=\sum_a\phi_aS_a,\quad b_r=x_r+iy_r.
\]
These are coordinates on the original spaces, not a rescaling of their
metrics. The mass metric read from \eqref{qh:hq-L} is
\[
 M=\operatorname{diag}\left(
   \frac{2\mathcal V_3}{g^2}I_{24},\
   \mathcal V_3I_8,\
   2\mathcal V_3I_6\right).
\]
We use its constant volume measure
\[
 d\mu_M=C_M\,d^{24}a\,d^8\phi\,d^3x\,d^3y,\qquad
 C_M=\sqrt{\det M}=\frac{2^{15}\mathcal V_3^{19}}{(g^2)^{12}},
\]
retaining this constant rather than setting it to one. The auxiliary
Hilbert space is \(\mathscr H=L^2(\mathcal Q,d\mu_M;\mathbb C)\).
In these coordinates canonical momentum is \(-i\hbar\partial\).
Thus the exact kinetic operator is
\begin{equation}\label{qh:hq-kinetic}
 T=-\frac{\hbar^2g^2}{4\mathcal V_3}\Delta_a
   -\frac{\hbar^2}{2\mathcal V_3}\Delta_\phi
   -\frac{\hbar^2}{4\mathcal V_3}(\Delta_x+\Delta_y).
\end{equation}
Equivalently the complex momentum is
\(\widehat\pi_{b_r}=-i\hbar\partial_{\bar b_r}
=-i\hbar(\partial_{x_r}+i\partial_{y_r})/2\);
its adjoint product is
\(-\hbar^2(\partial_{x_r}^2+\partial_{y_r}^2)/4\).
The Schr\"odinger equation is \(i\hbar\partial_t\psi=H\psi\).

% Root proof: exact scalar bound and combination with the matrix estimate.
\subsection{The retained cubic and a global scalar bound}
\label{qh:scalar-bound-section}

\begin{lemma}[Sharp cubic bound in the original complex coordinates]
\label{qh:scalar-bound}
For every traceless Hermitian $\Phi\in\operatorname{Herm}_3(\mathbb C)$
and $b\in\mathbb C^3$, put
\[
 s=\operatorname{tr}\Phi^2+2b^*b,\qquad
 N=\det\Phi-b^*\Phi b.
\]
Then
\begin{equation}\label{qh:cubic-bound}
 |N|\leq\frac{s^{3/2}}{3\sqrt6}.
\end{equation}
Consequently, with all original coefficients retained,
\begin{equation}\label{qh:scalar-coercivity}
 V(\Phi,b)=\frac{m^2s}{2}+\frac{\Gamma s^2}{4}-\kappa N
 \ \geq\ \frac{\Gamma s^2}{8}-C_\kappa,
 \qquad C_\kappa=\frac{\kappa^4}{54\Gamma^3},
\end{equation}
whenever $m^2\geq0$, $\Gamma>0$, and $\kappa\in\mathbb R$.
This is an inequality for the full potential, not its replacement by
the right-hand side.
\end{lemma}
\begin{proof}
Choose a unitary diagonalization of $\Phi$ and apply the same unitary
to $b$. The inverse conjugation restores both original fields; it
preserves $s,N$ and introduces no restriction on either field.
Write the three real eigenvalues as $\lambda_1,\lambda_2,\lambda_3$,
with sum zero, and let $z_i=|b_i|^2$, $u=\sum_i z_i$.
If $u>0$, the original cubic is the convex combination
\[
 N=\sum_{i=1}^3\frac{z_i}{u}
       \bigl(\lambda_1\lambda_2\lambda_3-u\lambda_i\bigr).
\]
For one index $i$ and the other two indices $j,k$, put
$d_i=\lambda_j-\lambda_k$. Trace zero gives
\[
 \lambda_j\lambda_k=(\lambda_i^2-d_i^2)/4,\qquad
 s=(3\lambda_i^2+d_i^2)/2+2u.
\]
Hence each term in parentheses is exactly
\[
 \lambda_i(\lambda_j\lambda_k-u)
       =\lambda_i^3-\frac{s\lambda_i}{2},\qquad
 |\lambda_i|\leq\sqrt{2s/3}.
\]
If $u=0$, the same identity equals $N$ for each $i$.
If $s=0$, both fields vanish and the assertion is immediate.
For $s>0$, the critical points of $x^3-sx/2$ on
$[-\sqrt{2s/3},\sqrt{2s/3}]$ are $x=\pm\sqrt{s/6}$.
At these two points and the two endpoints, its absolute value is
$s^{3/2}/(3\sqrt6)$. The elementary extreme-value theorem proves
the bound for every term and hence for their convex combination.
Equality is attained, for example, at
$\Phi=\operatorname{diag}(2r,-r,-r)$, $b=0$, $r\ne0$,
since then $s=6r^2$ and $N=2r^3$.

Put $\alpha=|\kappa|/(3\sqrt6)$ and $t=\sqrt{s}\geq0$.
Keep $m^2s/2$ as a nonnegative term when bounding the original
potential. The remaining estimate is
\[
 V\geq\frac{\Gamma t^4}{8}
       +\left(\frac{\Gamma t^4}{8}-\alpha t^3\right).
\]
For $\alpha>0$ the expression in parentheses has its global minimum
at $t=6\alpha/\Gamma$, with value
$-54\alpha^4/\Gamma^3=-\kappa^4/(54\Gamma^3)$.
For $\alpha=0$ it is nonnegative. This proves
\eqref{qh:scalar-coercivity} including both signs of $\kappa$.
\end{proof}


\begin{theorem}[Friedrichs realization with an explicit domain]
\label{qh:hq-Friedrichs}
Let \(W(q)=\mathcal V_3U(q)\) and
\[
 D=\operatorname{diag}\left(
   \frac{\hbar^2g^2}{4\mathcal V_3}I_{24},\
   \frac{\hbar^2}{2\mathcal V_3}I_8,\
   \frac{\hbar^2}{4\mathcal V_3}I_6\right),\qquad
 C=\frac{\mathcal V_3\kappa^4}{54\Gamma^3}.
\]
The differential expression \(T+W=-\nabla\cdot D\nabla+W\) on
\(C_c^\infty(\mathcal Q)\) has the following well-defined semibounded
Friedrichs realization \(H\). Its closed form and form domain are
\begin{align}
 \mathfrak h[\psi]&=\int\left((\nabla\psi)^*D\nabla\psi
                         +W|\psi|^2\right)d\mu_M,\notag\\
 \mathscr D&=H^1(\mathcal Q,d\mu_M)
                \cap L^2(\mathcal Q,(W+C)d\mu_M),\label{qh:hq-form}\\
 D(H)&=\{u\in\mathscr D:
          -\nabla\cdot D\nabla u+Wu\in L^2(\mathcal Q,d\mu_M)
                         \text{ in distributions}\}.\notag
\end{align}
The operator is self-adjoint, bounded below by \(-C\), and
\(C_c^\infty\) is a form core. No compactness assertion is needed
for this construction.
\end{theorem}
\begin{proof}
Lemma~\ref{qh:scalar-bound} proves directly in the original complex
coordinates that $V\geq\Gamma s^2/8-\kappa^4/(54\Gamma^3)$.
All the other terms of \(U\) are nonnegative. In particular \(W+C\)
is a continuous nonnegative polynomial, locally bounded everywhere.

The norm
\[
 \|u\|_{\mathscr D}^2
  =\int\left((\nabla u)^*D\nabla u+(W+C+1)|u|^2\right)d\mu_M
\]
is complete. Indeed a Cauchy sequence converges in \(H^1\) and in the
weighted \(L^2\) space, and their limits agree on every compact set.
This proves closedness. Smooth radial cutoffs converging to one
approximate any \(u\in\mathscr D\) in this norm: the potential and
gradient tails tend to zero, and the cutoff-gradient term is bounded
by a constant times \(R^{-2}\|u\|^2\). On a fixed compact set ordinary
mollification converges in \(H^1\), and the bounded weight there gives
weighted convergence. Thus \(C_c^\infty\) is a form core and the
form is densely defined.

For clarity one may construct its self-adjoint operator directly.
The positive form \(\mathfrak h+(C+1)\langle\cdot,\cdot\rangle\)
is an inner product on \(\mathscr D\). Riesz representation assigns
to every \(f\in\mathscr H\) a unique \(u=Rf\in\mathscr D\) such that
\[
 \mathfrak h(u,v)+(C+1)\langle u,v\rangle=\langle f,v\rangle
                  \quad(v\in\mathscr D).
\]
It follows that \(R\), viewed on \(\mathscr H\), is bounded,
self-adjoint, positive and injective; its range is dense because
its kernel is zero. The spectral theorem for bounded self-adjoint
operators gives the self-adjoint inverse \(R^{-1}\) on
\(\operatorname{Ran}R\). Set \(H=R^{-1}-(C+1)\).
This is precisely the Friedrichs operator of the original form core.
Integration by parts on compact test functions identifies its
distributional expression. Conversely, for \(u\in\mathscr D\) whose
displayed distribution is \(f\in L^2\), that integration-by-parts
identity extends from the form core to all of \(\mathscr D\);
the Riesz characterization gives \(u\in D(H)\), \(Hu=f\).
This proves exactly the asserted domain. It does not replace the
domain by an unjustified separate requirement \(Wu,\Delta u\in L^2\).
\end{proof}

\begin{theorem}[Spatially constant Gauss constraints and averaging]
\label{qh:hq-physical}
For \(g\in SU(3)\), let
\[
 R_gq=(gA_ig^{-1},g\Phi g^*,gb),\qquad
 (\mathcal U(g)\psi)(q)=\psi(R_{g^{-1}}q).
\]
This is a strongly continuous unitary representation on \(\mathscr H\).
For each \(\xi\in\mathfrak{su}(3)\), the operator
\[
 \widehat{\mathcal M}_\xi=-i\hbar K_\xi\cdot\nabla
                          \quad\hbox{on }C_c^\infty(\mathcal Q)
\]
is essentially self-adjoint. Its closure has domain
\[
 D(\widehat{\mathcal M}_\xi)
  =\{\psi\in\mathscr H:K_\xi\cdot\nabla\psi\in\mathscr H
                                      \text{ distributionally}\},
 \qquad
 \mathcal U(e^{t\xi})=e^{-it\widehat{\mathcal M}_\xi/\hbar}.
\]
The spatially constant physical Hilbert space and projection are
\begin{align}
 \mathscr H_{\rm phys}
  &=\{\psi:\mathcal U(g)\psi=\psi\ \forall g\in SU(3)\}
    =\bigcap_{a=1}^8\ker\widehat{\mathcal M}_{T_a},\notag\\
 \mathsf P\psi&=\frac1{\mathcal V_\mu}
       \int_{SU(3)}\mathcal U(g)\psi\,d\mu(g),\qquad
 \mathcal V_\mu=\int_{SU(3)}d\mu(g)>0 .\label{qh:hq-average}
\end{align}
Here the Haar mass is retained explicitly. The space is infinite
dimensional. The set \(C_c^\infty(\mathcal Q)^{SU(3)}
=\mathsf P C_c^\infty(\mathcal Q)\) is dense in it and is a form
core for its physical Hamiltonian
\[
 H_{\rm phys}=H|_{\mathscr H_{\rm phys}},\qquad
 D(H_{\rm phys})=D(H)\cap\mathscr H_{\rm phys}.
\]
Thus the physical Hamiltonian is self-adjoint and is exactly the
Friedrichs operator of the restricted form; no quotient-coordinate
measure or singular-orbit boundary condition has been inserted.
\end{theorem}
\begin{proof}
Each \(R_g\) preserves the three original inner products, hence the
mass metric, its measure and the kinetic form. It preserves \(U\)
by conjugation invariance of the trace/determinant and unitarity of
the triplet action. Continuity holds first on compact smooth
functions and then on \(L^2\) by density and unitarity.
The vector field \(K_\xi\) is linear with zero divergence: its flow
is an orthogonal linear map for each block. Its action on compact
smooth functions is therefore skew-adjoint under integration, so
\(-i\hbar K_\xi\cdot\nabla\) is symmetric. Differentiating the
displayed unitary action gives the stated generator sign.
Stone's theorem gives a self-adjoint generator. Its domain is
the displayed distributional domain: one direction follows by
differentiation against test functions; for the converse,
distributional integration along the flow gives
\[
 \mathcal U(e^{t\xi})\psi-\psi
   =-\int_0^t\mathcal U(e^{s\xi})
                   (K_\xi\cdot\nabla\psi)\,ds
\]
as an \(L^2\) identity, and strong continuity gives differentiability.

To verify that the original compact smooth domain is a core, use
smooth cutoffs depending on the blockwise invariant Euclidean radius.
They commute with \(K_\xi\). Convolution with a radial smooth kernel
also commutes with its orthogonal flow. Cutoff followed by convolution
therefore approximates both \(\psi\) and \(K_\xi\cdot\nabla\psi\)
in \(L^2\). This proves graph-core density and essential
self-adjointness. On that common smooth core,
\[
 [K_\xi\cdot\nabla,K_\eta\cdot\nabla]
   =-K_{[\xi,\eta]}\cdot\nabla,\qquad
 [\widehat{\mathcal M}_\xi,\widehat{\mathcal M}_\eta]
   =i\hbar\widehat{\mathcal M}_{[\xi,\eta]}.
\]
There is no divergence correction or ordering anomaly.

An invariant vector lies in every generator domain with derivative
zero. Conversely the eight vanishing distributional derivatives
give the vanishing derivative for every \(\xi\), hence invariance
under every one-parameter subgroup. These generate connected \(SU(3)\).
For example each unitary determinant-one matrix is unitarily
diagonalizable with eigenphases whose sum can be chosen to be zero,
so it is the exponential of a trace-free anti-Hermitian matrix.
This proves the kernel characterization with its domain included.

The Bochner integral in \eqref{qh:hq-average} is bounded in norm
by one. Haar invariance and inversion imply
\(\mathsf P^2=\mathsf P=\mathsf P^*\), with range precisely the
invariant vectors. Averaging a compact smooth function remains
smooth and has compact support, because the union of its supports
under a compact group is compact. Applying \(\mathsf P\) to any
compact-smooth approximation proves \(L^2\) density.
Nonzero smooth functions supported in pairwise disjoint annuli of
the invariant squared radius are invariant and orthogonal.
This proves infinite dimensionality by explicit test states.

The form and its domain are \(\mathcal U(g)\)-invariant.
The Riesz characterization above therefore gives
\(\mathcal U(g)R=R\mathcal U(g)\); averaging gives
\(\mathsf PR=R\mathsf P\). Functional calculus proves that
\(\mathsf P\) reduces \(H\), so its restriction has the claimed
self-adjoint domain. Jensen's inequality for the positive form
\(\mathfrak h+(C+1)\|\cdot\|^2\) makes \(\mathsf P\) a contraction
in the form norm. Averaging a form-core approximation of an
invariant vector proves the final form-core statement.
\end{proof}

% Integration copy with qh: labels and a_g,b_g; original proof preserved.
% Independent bounded pure-gauge matrix audit; no canonical edits.
\subsection*{An explicit quantum bound for the three retained gauge matrices}

Let $\mathfrak g=\mathfrak{su}(3)$ as a real eight-dimensional vector
space, with the retained metric
\[
 \langle X,Y\rangle=-\operatorname{tr}(XY),\qquad
 \|X\|^2=-\operatorname{tr}(X^2).
\]
For anti-Hermitian $X,Y$, the trace of $XY$ is real, by taking its
conjugate and using cyclicity. The metric is positive definite because
$-\operatorname{tr}(X^2)=\operatorname{tr}(X^*X)$.
Lebesgue measure and all Laplacians below are those of this metric,
not of the negative Killing form. Write $A=(A_1,A_2,A_3)\in\mathfrak g^3$.
For two arbitrary retained numbers $a_g,b_g>0$, consider on
$L^2(\mathfrak g^3,dA)$ the nonnegative quadratic form
\[
 q_{a_g,b_g}[\psi]=
 a_g\sum_{i=1}^3\int\|\nabla_{A_i}\psi\|^2\,dA+
 b_g\sum_{1\leq i<j\leq3}
       \int\|[A_i,A_j]\|^2|\psi|^2\,dA,                       \tag{MQ.1}
\]
initially on $C_c^\infty(\mathfrak g^3;\mathbb C)$.
The coefficients $a_g,b_g$ are independent retained positive numbers;
their subscripts distinguish them from the scalar data. No assertion about the scalar
coupling or a Gauss-reduced operator is included in the theorem below.

\begin{lemma}[The exact adjoint frequency sum, including root walls]
\label{qh:matrix-root-sum}
Let $X\in\mathfrak g$, and let $\theta_1\leq\theta_2\leq\theta_3$
be the real eigenvalues of $-iX$, with $\sum_p\theta_p=0$.
The positive semidefinite real operator
$K_X=-\operatorname{ad}_X^2$ has its eight eigenvalues, with real
multiplicity, equal to
\[
 0,0,\quad
 (\theta_1-\theta_2)^2,(\theta_1-\theta_2)^2,\quad
 (\theta_1-\theta_3)^2,(\theta_1-\theta_3)^2,\quad
 (\theta_2-\theta_3)^2,(\theta_2-\theta_3)^2.                  \tag{MQ.2}
\]
In particular, zeros in this list remain zeros at repeated eigenvalues.
Define the frequency sum with real multiplicities by
$\mathcal S(X)=\operatorname{tr}_{\mathbb R\mathfrak g}\sqrt{K_X}$.
Then
\[
 \mathcal S(X)
 =2\sum_{p<q}|\theta_p-\theta_q|
 =4(\theta_3-\theta_1),\qquad
 2\sqrt6\,\|X\|\leq\mathcal S(X)\leq4\sqrt2\,\|X\|.           \tag{MQ.3}
\]
Both constants are sharp. The lower equality at nonzero $X$ occurs
exactly when two eigenvalues coincide; the upper equality occurs exactly
when the eigenvalues are $(-r,0,r)$ for some $r\geq0$.
If each positive root is counted only once instead of using real
multiplicity, the sum is $\mathcal S(X)/2$ and its sharp lower
constant is $\sqrt6$.
\end{lemma}
\begin{proof}
Cyclicity gives
$\langle[X,Y],Z\rangle=-\langle Y,[X,Z]\rangle$.
Thus $\operatorname{ad}_X$ is skew-adjoint and
$\langle Y,K_XY\rangle=\|[X,Y]\|^2$.
Unitary diagonalization of $X$ is an isometry for the retained metric.
Multiplication of that unitary by a scalar phase adjusts its determinant
to one without changing its conjugation, so this diagonalization can
also be performed by $SU(3)$.

For diagonal $X=i\operatorname{diag}(\theta_1,\theta_2,\theta_3)$,
use the orthonormal Cartan vectors
\[
 D_1=\frac{i}{\sqrt2}\operatorname{diag}(1,-1,0),\qquad
 D_2=\frac{i}{\sqrt6}\operatorname{diag}(1,1,-2),
\]
and, for $p<q$, the orthonormal pair
\[
 U_{pq}=\frac{E_{pq}-E_{qp}}{\sqrt2},\qquad
 V_{pq}=\frac{i(E_{pq}+E_{qp})}{\sqrt2}.
\]
Direct multiplication gives
$[X,U_{pq}]=(\theta_p-\theta_q)V_{pq}$ and
$[X,V_{pq}]=-(\theta_p-\theta_q)U_{pq}$, while $D_1,D_2$
commute with $X$. These eight orthonormal vectors span $\mathfrak g$,
proving (MQ.2) and the exact sum in (MQ.3).

Put $r=\theta_2-\theta_1\geq0$ and $s=\theta_3-\theta_2\geq0$.
The original eigenvalues are exactly
\[
 \theta_1=-(2r+s)/3,\quad \theta_2=(r-s)/3,\quad
 \theta_3=(r+2s)/3,\quad
 \|X\|^2=\frac23(r^2+rs+s^2),\quad \mathcal S(X)=4(r+s).
\]
Consequently
\[
 \mathcal S(X)^2-24\|X\|^2=16rs,\qquad
 32\|X\|^2-\mathcal S(X)^2=\frac{16}{3}(r-s)^2.
\]
Their nonnegativity proves the bounds and their vanishing proves every
equality case. In particular the lower bound does not fail at a repeated
eigenvalue: at a nonzero wall exactly four real frequencies remain
positive. At $X=0$ both sides vanish.
\end{proof}

\begin{lemma}[Slice oscillators with all coefficients]
\label{qh:matrix-slice}
For fixed $X\in\mathfrak g$ and $\alpha,\beta>0$,
\[
 -\alpha\Delta_Y+\beta\|[X,Y]\|^2
 \ \geq\ \sqrt{\alpha\beta}\,\mathcal S(X)
 \quad\hbox{on }L^2(\mathfrak g,dY)                           \tag{MQ.4}
\]
as a quadratic-form inequality. No positive frequency is assigned to
the commuting coordinates.
\end{lemma}
\begin{proof}
For $\omega\geq0$ and $f\in C_c^\infty(\mathbb R;\mathbb C)$,
integration by parts proves the exact identity
\[
 \int|\sqrt\alpha f'+\sqrt\beta\,\omega y f|^2\,dy
 =\alpha\int|f'|^2\,dy+
   \beta\omega^2\int y^2|f|^2\,dy
   -\sqrt{\alpha\beta}\,\omega\int|f|^2\,dy .
\]
Indeed $2\operatorname{Re}(f'\bar f)=(|f|^2)'$ and the boundary
term vanishes. Thus each real oscillator contributes
$\sqrt{\alpha\beta}\,\omega$, including the correct zero contribution
when $\omega=0$.

Diagonalize the real symmetric $K_X$ in an orthonormal basis. This
orthogonal coordinate change preserves measure and $\Delta_Y$, and
the potential becomes the sum $\sum_{\nu=1}^8\omega_\nu^2 y_\nu^2$
with the frequencies in (MQ.2). Apply the preceding identity in
each coordinate, integrate the other coordinates, and sum to obtain
(MQ.4). This proof is for each fixed $X$. It requires no differentiable
choice of the diagonalizing basis as $X$ varies, and introduces no
derivatives in $X$.
\end{proof}

\begin{theorem}[An explicit coercive form bound and compact resolvent]
\label{qh:matrix-compact}
The closure of (MQ.1) has domain
\[
 \mathcal D(q_{a_g,b_g})=
 \left\{\psi\in H^1(\mathfrak g^3):
  \sum_{i<j}\int\|[A_i,A_j]\|^2|\psi|^2\,dA<\infty\right\}.    \tag{MQ.5}
\]
Let $H_A$ be its nonnegative Friedrichs operator. Its formal expression
on compactly supported smooth functions is exactly
\[
 H_A=-a_g\sum_i\Delta_{A_i}
                 +b_g\sum_{i<j}\|[A_i,A_j]\|^2.                \tag{MQ.6}
\]
For every $0\leq\vartheta<1$ the following form bound holds:
\[
 q_{a_g,b_g}[\psi]\geq
 \vartheta a_g\int\|\nabla\psi\|^2\,dA+
 2\sqrt{6(1-\vartheta)a_gb_g}
       \int\left(\sum_i\|A_i\|\right)|\psi|^2\,dA.             \tag{MQ.7}
\]
In particular, retaining half of the original kinetic energy gives
\[
 H_A\geq-\frac{a_g}{2}\Delta_{\mathfrak g^3}
                 +2\sqrt{3a_gb_g}\sum_i\|A_i\|.             \tag{MQ.8}
\]
The form-domain inclusion into $L^2(\mathfrak g^3)$ is compact.
Consequently $(H_A+\mu I)^{-1}$ is compact for every $\mu>0$,
and $H_A$ has only isolated eigenvalues of finite multiplicity,
with no finite accumulation point.
\end{theorem}
\begin{proof}
The potential in (MQ.6) is a continuous nonnegative polynomial.
The right side of (MQ.5), equipped with
$\|\psi\|_2^2+q_{a_g,b_g}[\psi]$, is complete: an associated Cauchy
sequence converges in $H^1$ and also after multiplication by the
square root of the potential; an almost-everywhere convergent
subsequence identifies the second limit as that product.
It is densely defined since it contains $C_c^\infty$.
Smooth cutoffs equal to one on a ball of radius $R$, zero outside
the ball of radius $2R$, with derivative bounded by a constant
times $R^{-1}$, approximate each member in this norm as $R\to\infty$.
The derivative error tends to zero by the $H^1$ norm and the
$R^{-1}$ bound; the weighted error tends to zero by integrability.
For each cutoff, convolution with a compactly supported smooth
approximate identity converges in $H^1$, and also in weighted
$L^2$ because the potential is bounded on the common compact
support. This proves the stated core and closure.

For the lower bound start on this core and split the quadratic form
exactly as
\[
 q_{a_g,b_g}[\psi]=
 \vartheta a_g\|\nabla\psi\|_2^2+
 \sum_{\substack{i,j\in\{1,2,3\}\\i\ne j}}
 \left\{\frac{(1-\vartheta)a_g}{2}\|\nabla_{A_j}\psi\|_2^2
       +\frac{b_g}{2}\int\|[A_i,A_j]\|^2|\psi|^2\,dA\right\}.
                                                               \tag{MQ.9}
\]
Every derivative appears in two directed pairs, and every unordered
potential pair also appears twice. Hence no kinetic or potential
term has been counted more than its original coefficient.
Apply (MQ.4) inside each directed summand with $X=A_i$,
$Y=A_j$, $\alpha=(1-\vartheta)a_g/2$, $\beta=b_g/2$.
Fubini integration in $A_i$ and in the third matrix gives the lower
bound
\[
 \frac{\sqrt{(1-\vartheta)a_gb_g}}2
       \int\mathcal S(A_i)|\psi|^2\,dA
\]
for that summand. There are two partners $j$ for each fixed $i$.
Their sum and (MQ.3) prove (MQ.7), and $\vartheta=1/2$ gives
(MQ.8). Extension to (MQ.5) follows from core convergence and
Fatou's lemma for the nonnegative position weight.

Here is a direct compactness proof. Fix $0<\vartheta<1$ and put
$c_\vartheta=2\sqrt{6(1-\vartheta)a_gb_g}>0$.
In any one retained orthonormal basis of $\mathfrak g^3=\mathbb R^{24}$,
let $Q_R=[-R,R]^{24}$. Outside this cube,
$\sum_i\|A_i\|\geq(\sum_i\|A_i\|^2)^{1/2}\geq R$. Thus
\[
 \int_{\mathbb R^{24}\setminus Q_R}|\psi|^2
 \leq\frac{q_{a_g,b_g}[\psi]}{c_\vartheta R},\qquad
 \int|\nabla\psi|^2\leq\frac{q_{a_g,b_g}[\psi]}{\vartheta a_g}.
                                                               \tag{MQ.10}
\]
Partition $Q_R$ into finitely many cubes of side $h=2R/N$,
and let $\Pi_{R,N}\psi$ be its mean on each small cube and zero
outside $Q_R$. This has finite-dimensional range. On any cube
$Q$ of side $h$,
\[
 \int_Q|\psi-\psi_Q|^2\leq h^2\int_Q|\nabla\psi|^2.            \tag{MQ.11}
\]
For completeness, in one dimension the variance identity gives
$\int_0^h|f-f_{\rm av}|^2
=(2h)^{-1}\int_0^h\int_0^h|f(x)-f(y)|^2\,dx\,dy$.
Cauchy--Schwarz in the fundamental theorem of calculus bounds
this by $h^2\int_0^h|f'|^2$. In a product cube, let $P_j$
be averaging in the $j$th coordinate. The commuting orthogonal
projections give the orthogonal decomposition
\[
 I-\prod_{j=1}^{24}P_j
 =\sum_{j=1}^{24}\left(\prod_{k<j}P_k\right)(I-P_j).
 \]
Applying the one-dimensional estimate to each summand and using
that prior-coordinate averaging contracts the derivative norm
proves (MQ.11), first for smooth functions and then by $H^1$
density.

Equations (MQ.10)--(MQ.11) give the explicit approximation bound
\[
 \|\psi-\Pi_{R,N}\psi\|_2^2
 \leq q_{a_g,b_g}[\psi]
       \left(\frac{h^2}{\vartheta a_g}
                    +\frac1{c_\vartheta R}\right).           \tag{MQ.12}
\]
On a bounded form-domain set choose first $R$ large and then
$N$ large. The error is uniformly as small as prescribed,
while $\Pi_{R,N}$ maps that set to a bounded subset of a fixed
finite-dimensional space. Such a set has a finite net at any
positive radius, by subdividing a bounding coordinate box.
Thus the original set is totally bounded in $L^2$, proving
compactness of the form-domain inclusion without a matrix-model
discreteness theorem.

The closed form defines its Friedrichs operator by the usual
Riesz construction: for each $f\in L^2$ and $\mu>0$ there is a
unique $u$ with
$q_{a_g,b_g}(u,v)+\mu\langle u,v\rangle=\langle f,v\rangle$
for every form-domain $v$. This follows by representing the
bounded functional on the Hilbert space with that form inner
product; coercivity is supplied by $\mu\|v\|_2^2$.
The solution is $(H_A+\mu I)^{-1}f$ and obeys
$q_{a_g,b_g}[u]+\mu\|u\|_2^2\leq\|f\|_2^2/\mu$.
Hence this resolvent factors through a bounded set in the form
domain and its compact inclusion, so it is compact.
The spectral theorem for the resulting compact self-adjoint
resolvent gives finite multiplicities and eigenvalues tending
to zero for the resolvent, equivalently eigenvalues of $H_A$
tending to infinity. Since the resolvent is injective and its
range is dense in an infinite-dimensional space, it has
infinitely many such eigenvalues. This proves all assertions.
\end{proof}

The inequality is a quadratic-form inequality, not a pointwise
lower bound on the quartic potential. That potential still vanishes
on every commuting triple. The exact mechanism is (MQ.9):
localizing the other matrices transversely to a large fixed matrix
costs the nonzero oscillator ground energies in (MQ.4); their
sum remains coercive even at a root wall.

\begin{proposition}[An additional explicit lower bound on the energy origin]
\label{qh:matrix-bottom}
With precisely the energy origin in (MQ.6),
\[
 H_A\geq 9(3/2)^{1/3}a_g^{2/3}b_g^{1/3}\,I.                    \tag{MQ.13}
\]
This is a lower bound on the bottom of the operator, not a claimed
lower bound of the same size on the difference between its first two
energy levels.
\end{proposition}
\begin{proof}
First, for any $\alpha,\gamma>0$,
\[
 -\alpha\frac{d^2}{dx^2}+\gamma|x|
 \geq \frac3{2^{8/3}}(\alpha\gamma^2)^{1/3}.                 \tag{MQ.14}
\]
To prove it, take $f\in C_c^\infty(\mathbb R)$ with $\|f\|_2=1$,
and write $t=\|f'\|_2^2$ and $M=\int|x||f(x)|^2\,dx>0$.
Integration of $(|f|^2)'$ from both ends of the line gives
$\|f\|_\infty^2\leq\|f\|_2\|f'\|_2=\sqrt t$.
Markov's inequality puts at least half of the mass inside
$[-2M,2M]$, so
$1/2\leq4M\|f\|_\infty^2\leq4M\sqrt t$.
Thus $tM^2\geq1/64$. Minimizing
$\alpha/(64M^2)+\gamma M$ over $M>0$, whose unique minimum
occurs at $M^3=\alpha/(32\gamma)$, proves (MQ.14).
Homogeneity in $f$ removes the unit-norm restriction.

Write $x_{i\nu}$ for the original metric-orthonormal real coordinates,
$1\leq i\leq3$, $1\leq\nu\leq8$. Cauchy--Schwarz gives
$\|A_i\|\geq8^{-1/2}\sum_{\nu=1}^8|x_{i\nu}|$.
By (MQ.8),
\[
 H_A\geq
 \sum_{i,\nu}\left\{
    -\frac{a_g}{2}\partial_{x_{i\nu}}^2
           +\sqrt{\frac{3a_gb_g}{2}}\,|x_{i\nu}|\right\}.
\]
Apply (MQ.14) in each of these twenty-four coordinates and
integrate the others. Its total constant is exactly
\[
 24\frac3{2^{8/3}}\left(\frac{a_g}{2}\,\frac{3a_gb_g}{2}\right)^{1/3}
 =9(3/2)^{1/3}a_g^{2/3}b_g^{1/3}.
\]
Core closure gives the operator-form bound stated.
\end{proof}

\paragraph{Exact relation to the retained homogeneous gauge density.}
The canonical density uses signature $(+---)$ and anti-Hermitian
gauge fields, and its pure gauge part is
\[
 \mathcal L_{\rm gauge,hom}
 =\frac1{g^2}\sum_i\|E_i\|^2
       -\frac1{g^2}\sum_{i<j}\|[A_i,A_j]\|^2,\qquad
 E_i=\dot A_i+[A_0,A_i].
\]
The full scalar potential retained there is
$V=m^2s/2+\Gamma s^2/4-\kappa(\det\Phi-b^*\Phi b)$,
with $s=\operatorname{tr}\Phi^2+2b^*b$.
It is not part of the pure matrix operator (MQ.6).
For the pure gauge homogeneous kinetic form on a finite torus of
retained volume $\mathcal V_3>0$, fixing temporal gauge at the
unreduced coordinate level gives
$L=\mathcal V_3\sum_i\|\dot A_i\|^2/g^2
-\mathcal V_3\sum_{i<j}\|[A_i,A_j]\|^2/g^2$.
Its canonical momenta are exactly
$p_i=2\mathcal V_3\dot A_i/g^2$; the Legendre transform is
\[
 H_{\rm cl}=\frac{g^2}{4\mathcal V_3}\sum_i\|p_i\|^2
              +\frac{\mathcal V_3}{g^2}
                       \sum_{i<j}\|[A_i,A_j]\|^2.
\]
The specified canonical substitution $p_i=-i\hbar\nabla_{A_i}$
with $\hbar>0$ therefore gives the exact coefficient map
\[
 a_g=\frac{\hbar^2g^2}{4\mathcal V_3},\qquad
 b_g=\frac{\mathcal V_3}{g^2},\qquad
 a_gb_g=\frac{\hbar^2}{4}.                                  \tag{MQ.15}
\]
No coupling or volume is set to one. Under this map, (MQ.8) reads
\[
 H_A\geq-\frac{\hbar^2g^2}{8\mathcal V_3}\Delta
                   +\sqrt3\,\hbar\sum_i\|A_i\|,
\]
and (MQ.13) reads
$H_A\geq9(3\hbar^4g^2/(32\mathcal V_3))^{1/3}I$.
These coefficient identities concern the displayed unreduced
homogeneous quantum model. They impose neither the Gauss constraint
nor the omitted scalar coupling, and assert no continuum or
infinite-volume field-theoretic gap.

\paragraph{Primary literature and scope.}
The discreteness phenomenon for bosonic commutator matrix potentials
is classical. B.~Simon, \emph{Some quantum operators with discrete
spectrum but classically continuous spectrum}, Annals of Physics
\textbf{146} (1983), 209--220, DOI
\url{https://doi.org/10.1016/0003-4916(83)90057-X}, treats a
compact semisimple Lie-algebra matrix operator in equation~(3)
and Section~7, Corollary~4. Section~2 gives the elementary
zero-point-oscillator argument for the two-variable model.
The present proof combines a direct slice argument with the
explicit $\mathfrak{su}(3)$ root sum in the metric
$-\operatorname{tr}(XY)$ and proves compactness directly by
(MQ.12). It does not use the Fefferman--Phong discreteness theorem
employed in that paper's Section~7, nor claim the underlying
discreteness result as new.

The related finite-volume Yang--Mills analysis is M.~L\"uscher,
\emph{Some analytic results concerning the mass spectrum of Yang--Mills
gauge theories on a torus}, Nuclear Physics B \textbf{219} (1983),
233--261, DOI \url{https://doi.org/10.1016/0550-3213(83)90436-4}.
That field-theoretic small-volume setting is not substituted for the
precisely specified matrix operator here; none of its perturbative
conclusions or coefficients is assumed in the foregoing proof.


\subsection{Compactness for the entire coupled Hamiltonian}
\label{qh:combined-section}

\begin{theorem}[All thirty-eight coordinates are quantum confined]
\label{qh:full-compact}
Retain the configuration metric
\[
 \gamma_E=\operatorname{tr}(d\Phi^2)+2\,db^*db,
 \qquad E=\operatorname{Herm}_3(\mathbb C)_0\oplus\mathbb C^3,
\]
and the metric $-\operatorname{tr}$ on each of the three
$\mathfrak{su}(3)$ factors. Let $T_A=-\sum_i\Delta_{A_i}$ and
$T_E=-\Delta_{\gamma_E}$; these are nonnegative Laplacians.
For every $0<\eta<1$, the full Hamiltonian $H_{\rm hom}=H$ obeys
the quadratic-form inequality
\begin{align}
 H_{\rm hom}\ \geq\;&
 \frac{\eta\hbar^2g^2}{4\mathcal V_3}T_A
 +\frac{\hbar^2}{2\mathcal V_3}T_E
 +\hbar\sqrt{6(1-\eta)}\sum_{i=1}^3\|A_i\|
 \notag\\*[-1mm]
 &+\frac{\mathcal V_3\Gamma}{8}s^2
 -\mathcal V_3 C_\kappa.
 \label{qh:full-form-bound}
\end{align}
Its Friedrichs realization on the entire thirty-eight-dimensional
configuration space has compact resolvent. No commuting or scalar
directions have been deleted in this assertion.
\end{theorem}
\begin{proof}
Apply the matrix estimate with
$a_g=\hbar^2g^2/(4\mathcal V_3)$ and
$b_g=\mathcal V_3/g^2$. Their product is exactly $\hbar^2/4$.
The scalar-gradient and gauge--scalar potentials
\[
 \frac{\mathcal V_3}{2}\sum_i\operatorname{tr}[A_i,\Phi]^2,
 \qquad \mathcal V_3\sum_i\|A_i b\|^2
\]
are nonnegative: $[A_i,\Phi]$ is Hermitian. Retain these terms
in the definition of the operator and use their signs only in
this lower bound. Lemma~\ref{qh:scalar-bound} now proves
\eqref{qh:full-form-bound} on compactly supported smooth
functions, and core closure extends it to the form domain.

Write $q$ for the closed form and
$\widetilde q[\psi]=q[\psi]+\mathcal V_3C_\kappa\|\psi\|_2^2$.
The preceding bound makes $\widetilde q$ nonnegative.
In the original metric the squared radial coordinate is
$r^2=\sum_i\|A_i\|^2+s$.
The retained positive weight
\[
 W_\eta=\hbar\sqrt{6(1-\eta)}\sum_i\|A_i\|
                       +\mathcal V_3\Gamma s^2/8
\]
tends to infinity as $r\to\infty$. More explicitly, outside
the ball $r<R$ either $\sum_i\|A_i\|^2\geq R^2/2$ or
$s\geq R^2/2$, so
\[
 W_\eta\geq
 w_R:=\min\left\{\hbar\sqrt{3(1-\eta)}R,
                        \mathcal V_3\Gamma R^4/32\right\}.
\]
Both coefficients are strictly positive for the stated parameters.
Thus
\[
 \int_{r\geq R}|\psi|^2\leq\widetilde q[\psi]/w_R.
\]
The derivative bound from \eqref{qh:full-form-bound} is
\[
 c_0\int|\nabla\psi|^2\leq\widetilde q[\psi],\qquad
 c_0=\min\{\eta\hbar^2g^2/(4\mathcal V_3),
                          \hbar^2/(2\mathcal V_3)\}>0,
\]
where the gradient uses the product metric just stated.
Metric-orthonormal coordinates are an isometric representation,
not a change to any coupling, original field, or volume.
On the cube $[-R,R]^{38}$, replace a function by its averages on
finitely many cubes of side $h=2R/N$ and by zero outside the
large cube. The coordinate-averaging argument in the matrix proof
works in thirty-eight coordinates without changing its constant:
the squared local error is at most $h^2\int|\nabla\psi|^2$.
The cube's complement lies outside $r<R$, so the total error is at most
\[
 \widetilde q[\psi]\left(h^2/c_0+1/w_R\right).
\]
On every form-norm bounded set this tends uniformly to zero by
first increasing $R$ and then increasing $N$. The approximations
have finite-dimensional range, and their norms are bounded by
$\|\psi\|_2$. Therefore that set is totally bounded in $L^2$.
The complete Hilbert form domain has compact inclusion in $L^2$.
The Riesz construction with the positive form $\widetilde q+\|\cdot\|_2^2$ maps bounded
$L^2$ sets to bounded form-domain sets. Its composition with this
inclusion is the resolvent, proving compactness for the full operator.
\end{proof}

The original classical potential remains zero in its pure-gauge
part for every commuting triple. Inequality~\eqref{qh:full-form-bound}
uses the derivative cost in directions transverse to such triples.
The strictly positive coefficient in front of $\sum_i\|A_i\|$
does not assert a pointwise classical potential that was absent
from the original density.

\begin{theorem}[Strict positivity and the bottom spectral projection]
\label{qh:hq-positive}
For the actual polynomial Hamiltonian above and every \(\tau>0\),
\(e^{-\tau H}\) is positivity improving: if \(0\ne f\geq0\), then
\(e^{-\tau H}f>0\) almost everywhere. The same holds on the invariant
Hilbert space. Moreover
\[
 \inf\sigma(H_{\rm phys})=\inf\sigma(H)=:E_0,\qquad
 \operatorname{rank}\mathbf1_{\{E_0\}}(H)\leq1.
\]
Every nonzero vector in that bottom eigenspace is a scalar multiple
of one strictly positive normalized, \(SU(3)\)-invariant real
function. These conclusions hold without assuming compactness.
\end{theorem}
\begin{proof}
We give the heat-kernel argument with its constants. In the above real
coordinates the free kernel relative to Lebesgue measure is
\[
 p_\tau(q,r)=(4\pi\tau)^{-19}(\det D)^{-1/2}
        \exp\!\left(-\frac{(q-r)^{\mathsf T}D^{-1}(q-r)}{4\tau}\right).
\]
Let \(B_s\), \(0\leq s\leq\tau\), be a standard 38-dimensional
Brownian motion and form the bridge
\[
 Z_s=(1-s/\tau)q+(s/\tau)r
                   +\sqrt{2D}\,(B_s-(s/\tau)B_\tau).
\]
The Friedrichs kernel, relative to the retained measure \(d\mu_M\), is
\begin{equation}\label{qh:hq-FK}
 K_\tau(q,r)=\frac{p_\tau(q,r)}{C_M}\,
       \mathbb E\exp\!\left(-\int_0^\tau W(Z_s)\,ds\right).
\end{equation}
Here \(\tau\) is the semigroup parameter with inverse-energy units;
Euclidean physical time is \(\hbar\tau\). No value of \(\hbar\)
has been set to one in \(D\).

For completeness, the formula can be justified within the present
bounded-below polynomial setting. Put
\(W_n=\min(W+C,n)\). For bounded \(W_n\), the Duhamel iteration of
the free semigroup converges in norm, since its \(j\)-th ordered
term is bounded by \((n\tau)^j/j!\). The Markov Gaussian densities
of the bridge identify exactly these ordered integrals with the
power-series expansion of
\(\mathbb E\exp(-\int_0^\tau W_n(Z_s)\,ds)\).
Thus the kernel formula holds for the bounded potentials.
The forms with potentials \(W_n\) increase to the form with
potential \(W+C\). Their resolvents converge strongly: for
\(u_n=(T+W_n+\lambda)^{-1}f\), \(\lambda>0\), the variational
identities bound \(u_n\) in \(H^1\); testing each fixed truncated
form and passing weakly shows that any weak limit belongs to
the limiting form domain. Compactly supported tests identify
its unique limiting resolvent equation. The identities
\[
 \mathfrak h_n[u_n]+\lambda\|u_n\|^2
       =\operatorname{Re}\langle f,u_n\rangle
\]
and lower semicontinuity then force convergence of the \(L^2\)
norms, hence strong convergence. The spectral functional calculus
gives strong convergence of their semigroups.
Every bridge path is continuous on a compact interval and has
bounded image. The continuous polynomial \(W\) is therefore
bounded on that path, and monotone convergence of its truncated
integrals and dominated convergence of their exponentials yield
\eqref{qh:hq-FK}, restoring the factor \(e^{C\tau}\).
This is the bounded-below Brownian-bridge Feynman--Kac formula,
also stated in B.~Simon, \emph{A Feynman--Kac Formula for Unbounded
Semigroups}, arXiv:math-ph/9907022, Theorem~1.1,
\url{https://arxiv.org/abs/math-ph/9907022}.

In particular the exponential in \eqref{qh:hq-FK} is strictly
positive almost surely and bounded above by \(e^{C\tau}\).
Since \(p_\tau(q,r)>0\), the kernel is strictly positive for all
\(q,r\). Its Gaussian upper bound makes it a valid \(L^2\)
kernel action, and integration against any nonzero nonnegative
function proves positivity improvement. The invariant subspace
is reducing and the action is the same positive kernel, giving
its assertion as well.

To prove equality of the two spectral infima without presupposing
an eigenfunction, take any \(f\in\mathscr D\) and define
\[
 F(q)=\left(\frac1{\mathcal V_\mu}
       \int_{SU(3)}|f(R_{g^{-1}}q)|^2\,d\mu(g)\right)^{1/2}.
\]
This is an invariant nonnegative function with \(\|F\|=\|f\|\).
Apply the norm inequality to the Hilbert-valued Sobolev function
\(q\mapsto(g\mapsto f(R_{g^{-1}}q))\). More explicitly, differentiating
\((\int|f(R_{g^{-1}}q)|^2+\varepsilon^2)^{1/2}\), using
Cauchy--Schwarz, and then taking \(\varepsilon\downarrow0\), gives
\[
 \int(\nabla F)^*D\nabla F\,d\mu_M
 \leq \frac1{\mathcal V_\mu}\int_{SU(3)}
       \int(\nabla\mathcal U(g)f)^*D
                           \nabla\mathcal U(g)f\,d\mu_M\,d\mu(g).
\]
One first argues locally for smooth functions and then uses form
approximation and lower semicontinuity. Invariance of \(W\) gives
equality of the potential integrals. Thus \(F\in\mathscr D\) and
\(\mathfrak h[F]\leq\mathfrak h[f]\). The variational
characterization of the spectral infimum proves
\(\inf\sigma(H_{\rm phys})\leq\inf\sigma(H)\); the reverse inequality
follows from restriction to a subspace.

Finally suppose the bottom spectral projection is nonzero, and take
a real nonzero bottom eigenfunction \(u\); real and imaginary parts
are eigenfunctions because \(H\) is real. The weak-gradient inequality
\(|\nabla|u||\leq|\nabla u|\) gives
\(\mathfrak h[|u|]\leq E_0\|u\|^2\). The lower spectral bound forces
equality, and the spectral theorem then gives
\(|u|\in\ker(H-E_0)\). Positivity improvement makes \(|u|>0\)
almost everywhere. A sign-changing real bottom eigenfunction is
impossible: its positive and negative parts are both sent to
strictly positive functions, so
\[
 |e^{-\tau H}u|<e^{-\tau H}|u|
\]
almost everywhere, contradicting the eigenfunction identities.
Two independent real bottom eigenfunctions would yield, by
orthogonalization against the positive one, such a sign-changing
eigenfunction. Thus the complex eigenspace has dimension one.
Every \(\mathcal U(g)|u|\) is another positive normalized bottom
eigenfunction, so its one-dimensional proportionality constant
is positive and of modulus one, hence one. This proves gauge
invariance and all assertions of the theorem.
\end{proof}

\begin{corollary}[Unique physical ground state and a finite-model gap]
\label{qh:hq-ground}
The actual homogeneous quantum Hamiltonian has one normalized positive
ground state, unique up to phase and invariant under spatially constant
\(SU(3)\). The physical eigenvalues, counted with multiplicity, satisfy
\[
 E_0=E^{\rm phys}_0<E^{\rm phys}_1\leq E^{\rm phys}_2\leq\cdots,
 \qquad E^{\rm phys}_n\longrightarrow+\infty .
\]
In particular \(E^{\rm phys}_1-E_0>0\), without a numerical value being
claimed.
\end{corollary}
\begin{proof}
The positive compact resolvent $(H+C+1)^{-1}$ proved above has a nonzero top eigenvalue,
whose inverse supplies an eigenfunction at the bottom of \(H\).
The strict-positivity theorem makes it unique and invariant.
The physical space is infinite dimensional, and its compact
resolvent has finite-multiplicity eigenvalues accumulating only
at zero. Applying the inverse affine spectral relation proves
the displayed sequence and its strict first inequality.
\end{proof}
This is a gap of the specified finite-dimensional quantum mechanics.
No numerical value follows from positivity, and classical slow
solutions do not compute its level spacings. The proof retains fixed
positive \(\hbar,g^2,\mathcal V_3,\Gamma\); it asserts no uniform
gap in limits of these parameters or in a continuum limit.

\begin{proposition}[Exact classical embedding and its quotient limitations]
\label{qh:hq-embedding}
Let \(\iota\) send each homogeneous configuration and velocity to
the corresponding spatially constant fields on the trivial
\(SU(3)\) bundle over \(\mathbb T^3\).
The full action pulls back to \eqref{qh:hq-L}. Its homogeneous
Euler--Lagrange and Gauss solutions are exact full classical solutions.
On phase space use constant momentum densities
\(\Pi_i/\mathcal V_3,\Pi_\Phi/\mathcal V_3,\pi_b/\mathcal V_3\);
then the full canonical one-form pulls back to \(\vartheta\)
exactly. The map is equivariant for spatially constant \(SU(3)\).
The induced map from its constant-group orbit space to the full
torus gauge orbit space is not injective.
\end{proposition}
\begin{proof}
Spatial derivatives of the retained coefficients vanish. Substitution
in the density gives \eqref{qh:hq-L}, and the full local
Euler--Lagrange residuals are then constant in space.
Vanishing of their homogeneous variations for every time-dependent
matrix variation therefore makes those residuals zero pointwise,
including the \(A_0\) residual. They consequently vanish against
arbitrary, not necessarily homogeneous, spacetime test variations.
This proves the exact classical implication without a symmetry
reduction heuristic. Integrating constant momentum densities against
constant configuration variations multiplies them by \(\mathcal V_3\),
proving the asserted one-form pullback and thus its symplectic
pullback. Constant-group equivariance is direct.

For the limitation, let \(L_1\) be the first torus period and set
\[
 g(x_1)=\operatorname{diag}
   (e^{2\pi ix_1/L_1},e^{-2\pi ix_1/L_1},1).
\]
With \(D=\partial+A\), gauge transformation is
\[
 A_i'=gA_ig^{-1}-(\partial_i g)g^{-1},\quad
 \Phi'=g\Phi g^*,\quad b'=gb.
\]
The homogeneous state \(A_i=0,\Phi=b=0\) is thereby equivalent in
the full torus theory to the homogeneous state
\[
 A'_1=-\frac{2\pi i}{L_1}\operatorname{diag}(1,-1,0),\qquad
 A'_2=A'_3=0,\qquad \Phi'=b'=0 .
\]
No constant conjugation sends the zero \(A_1\) to this nonzero
matrix, proving noninjectivity of the orbit-space map.
On the other hand the same spatial gauge function sends the constant
\(\Phi=E_{12}+E_{21}\) to
\[
 e^{4\pi ix_1/L_1}E_{12}+e^{-4\pi ix_1/L_1}E_{21},
\]
which is not homogeneous. Thus spatial gauge maps do not act on
the whole homogeneous ansatz. More generally, a gauge map preserving
every homogeneous triplet \(b=e_j\), \(j=1,2,3\), must have all
three columns constant, and hence must itself be spatially constant.
This argument applies in particular to any genuinely large
transformation not homotopic to a constant. The explicit diagonal
example above is not being mislabeled as topologically large in
\(SU(3)\); its periodicity alone already proves the required
quotient obstruction.
\end{proof}

\paragraph{Why the classical embedding is not a quantum-sector theorem.}
There is a direct interaction obstruction already in a finite Fourier
cutoff. Choose the trace-free Hermitian
\(S=\operatorname{diag}(1,-1,0)/\sqrt2\), so
\(\operatorname{tr}S^2=1\), \(\det S=0\), and put
\[
 \Phi(x)=\bigl(q+y\cos(2\pi x_1/L_1)\bigr)S,\qquad A_i=b=0 .
\]
The quartic part of the spatially integrated potential is exactly
\[
 \frac{\Gamma\mathcal V_3}{4}
       \left(q^4+3q^2y^2+\frac38y^4\right).
\]
Indeed the averages of the cosine, its square, its cube and its
fourth power are \(0,1/2,0,3/8\). The mass and gradient terms
separate in \(q,y\), the cubic determinant vanishes, but the mixed
term \(3\Gamma\mathcal V_3q^2y^2/4\) remains.
The kinetic density is
\(\mathcal V_3\dot q^2/2+\mathcal V_3\dot y^2/4\);
thus this example also retains the unequal Fourier-mode masses.

A putative product sector consisting of all \(q\)-wavefunctions
times one fixed nonzero normalizable \(y\)-wavefunction \(\Omega\)
is not invariant under this interacting cutoff Hamiltonian.
Here invariance means a reducing quantum subspace, equivalently that
its orthogonal projection commutes with the unitary evolution.
This includes the necessary domain issue: if that projection
commuted with the Hamiltonian, applying it to
\(f(q)v(y)\in C_c^\infty\), with \(\langle\Omega,v\rangle\ne0\),
would put \(f(q)\Omega(y)\) in the operator domain for every
\(f\in C_c^\infty(\mathbb R)\). Its distributional Hamiltonian is
the separated \(q\)-term plus
\(f(q)H_y\Omega+cq^2f(q)y^2\Omega\),
\(c=3\Gamma\mathcal V_3/4\). Integrating against two compact
\(q\)-test functions with distinct averaged \(q^2\) values
isolates \(H_y\Omega\) and \(y^2\Omega\) as \(L^2\) functions.
Thus the following projections are legitimate even if no regularity
of \(\Omega\) was assumed initially.
Projecting its action orthogonally to \(\Omega\), for arbitrary
compact smooth \(q\)-wavefunctions, would require
\[
 (H_y\Omega)_\perp+
       \frac{3\Gamma\mathcal V_3}{4}q^2(y^2\Omega)_\perp=0
                         \quad\hbox{for almost every }q .
\]
Its two coefficients must vanish, implying
\(y^2\Omega=\lambda_\Omega\Omega\). Multiplication by \(y^2\) has no nonzero
\(L^2(\mathbb R)\) eigenfunction: its level sets have Lebesgue
measure zero. Thus no fixed discarded-mode vacuum yields that
product invariant sector. The distribution \(\delta(y)\), which
would impose a literally vanishing discarded coordinate, is not
an \(L^2\) state either.
This does not rule out every possible correspondence with a full
quantum theory; it disproves the specific automatic-subsector
inference and states exactly which finite model has been quantized.
No continuum quantum Hilbert space, full local-gauge quotient,
renormalized Hamiltonian, or continuum mass-gap value is constructed
by the finite-dimensional argument.

\paragraph{Gauss-law reference.}
D.~Tong, \emph{Lectures on Gauge Theory}, Section 2.2.1,
equations (2.27)--\allowbreak(2.31),
\url{https://www.damtp.cam.ac.uk/user/tong/gaugetheory.html}, retains the Gauss equation after temporal
gauge and imposes it on physical states. His Hermitian-field
normalization is not substituted for the retained anti-Hermitian
normalization here; the exact factors and the sign of
\eqref{qh:hq-Gauss} were derived above from the given density.
For this compact spatial torus, spatially constant \(SU(3)\) is
the finite gauge subgroup being imposed, rather than an asymptotic
charge group at spatial infinity.
\par
\section{An exact moving-ground-state comparison for the retained Fourier pair}
\label{fq:section}

The result below concerns precisely the two-coordinate variational
truncation
\[
 \Phi(x)=\bigl(q+y\cos(kx_1)\bigr)S,\qquad
 S=\operatorname{diag}(1,-1,0)/\sqrt2,\qquad
 k=2\pi/L_1,\qquad A_i=b=0.
\]
It does not assume that this two-dimensional plane is an invariant
subspace of the full classical field equations or a reducing subspace
of an independently constructed full quantum field theory.
Throughout, retain
\[
 \hbar>0,\quad\mathcal V_3>0,\quad\Gamma>0,\quad L_1>0,
 \qquad m^2\geq0,
\]
\[
 \qquad \omega_*:=\sqrt{m^2+k^2}>0,
 \qquad \omega(q):=\sqrt{\omega_*^2+3\Gamma q^2}.
\]
All bounds will refer to these original parameters and coordinates.
The quartic discarded-mode interaction is never removed from the operator.

\subsection{Exact coefficients, measures, and self-adjoint operators}

Since $\operatorname{tr}S^2=1$ and $\det S=0$, the pullback of the
cubic determinant is zero. The averages of $\cos(kx_1)$ and its first
four powers are $0,1/2,0,3/8$ for powers one through four. The exact
integrated Lagrangian is therefore
\begin{align}
 L_2={}&\frac{\mathcal V_3}{2}\dot q^2
       +\frac{\mathcal V_3}{4}\dot y^2
       -V_q(q)-\frac{\mathcal V_3\omega(q)^2}{4}y^2-cy^4,
       \label{fq:L}\\
 V_q(q):={}&\mathcal V_3\left(\frac{m^2q^2}{2}
                         +\frac{\Gamma q^4}{4}\right),\qquad
 c:=\frac{3\Gamma\mathcal V_3}{32}>0.\notag
\end{align}
In particular the gradient term is $\mathcal V_3k^2y^2/4$ and the
mixed quartic term is $3\Gamma\mathcal V_3q^2y^2/4$.
The momenta are $p_q=\mathcal V_3\dot q$ and
$p_y=\mathcal V_3\dot y/2$. Thus the mass metric and its measures are
\[
 M_2=\operatorname{diag}(\mathcal V_3,\mathcal V_3/2),\quad
 d\mu_q=\sqrt{\mathcal V_3}\,dq,\quad
 d\mu_y=\sqrt{\mathcal V_3/2}\,dy,\quad
 d\mu_2=\frac{\mathcal V_3}{\sqrt2}\,dq\,dy=d\mu_qd\mu_y.
\]
We retain these constants. Put
\[
 \mathscr K=L^2(\mathbb R,d\mu_y),\qquad
 \mathscr L=L^2(\mathbb R,d\mu_q),\qquad
 \mathscr H_2=L^2(\mathbb R,d\mu_q;\mathscr K).
\]
\[
 a:=\frac{\hbar^2}{2\mathcal V_3},\quad
 a_y:=\frac{\hbar^2}{\mathcal V_3},\quad
 b_q:=\frac{\mathcal V_3\omega(q)^2}{4}.
\]
The Hamiltonian is exactly
\begin{equation}\label{fq:H}
 H_2=-a\partial_q^2+V_q(q)+h(q),\qquad
 h(q)=-a_y\partial_y^2+b_qy^2+cy^4.
\end{equation}
The symbols $a,a_y,b_q,c$ denote these retained coefficients, not a
rescaling of either variable.

Every displayed nonnegative Schr\"odinger expression is realized by
its Friedrichs form. For example
\[
 \mathcal D_y=H^1(\mathbb R,d\mu_y)
                  \cap\{u:y^2u\in\mathscr K\},\qquad
 h_q[u]=a_y\|u'\|^2+b_q\|yu\|^2+c\|y^2u\|^2.
\]
Its form domain is independent of $q$, since $b_q,c>0$.
The full form domain consists of $\Psi\in H^1(\mathbb R^2,d\mu_2)$
for which the original nonnegative multiplication potential in
\eqref{fq:H} has finite integral against $|\Psi|^2$.
Completeness follows by convergence in $H^1$ and after multiplication
by the square root of that potential; an almost-everywhere subsequence
identifies the limits. Smooth cutoffs and then compactly supported
mollification prove that $C_c^\infty$ is a form core. The operator
domain is the form-domain vectors whose displayed differential
expression belongs to $L^2$ distributionally, as follows by testing
against the core and using the Riesz representation of the shifted form.
Since $b_qy^2+cy^4\to\infty$ in $y$, the fibre form-domain inclusion
is compact: the potential controls the $L^2$ tail and the derivative
controls finite-cell averaging errors on bounded intervals.
The same argument proves compactness for $H_2$, using
$V_q(q)+cy^4\to\infty$ as $q^2+y^2\to\infty$.

We shall also use the following elementary operator-core fact.
For a smooth nonnegative real potential on $\mathbb R^n$ and a
constant positive kinetic matrix, $C_c^\infty$ is an operator core
for the Friedrichs operator. Indeed an $L^2$ distributional solution
of $(H^*+1)u=0$ is locally $H^2$ by local elliptic regularity.
Testing against $\chi_R^2u$, with a smooth compact cutoff whose
gradient is $O(R^{-1})$, gives
\[
 \int\bigl(|\nabla(\chi_Ru)|_D^2+(V+1)|\chi_Ru|^2\bigr)
 =\int|\nabla\chi_R|_D^2|u|^2=O(R^{-2})\|u\|^2.
\]
Letting $R\to\infty$ proves $u=0$. Thus the range of $H+1$ on
the minimal domain is dense; its lower bound makes the range of
its closure closed. That closure is surjective at $-1$ and symmetric,
hence self-adjoint. It is therefore the Friedrichs realization and
the stated domain is an operator core. This argument applies to
the fibres, the full operator, and the effective operator below.

\subsection{The actual normalized positive fibre ground state}

\begin{lemma}[Fibre energies and a proved quantitative separation]
\label{fq:gap}
For every $q$, $h(q)$ has a simple lowest eigenvalue $e(q)$ and a
unique real strictly positive eigenfunction $\Omega_q$ with
$\|\Omega_q\|_{\mathscr K}=1$. Write $e_1(q)$ for its second
eigenvalue, counting multiplicity. Then
\begin{equation}\label{fq:energy-bounds}
 \frac{\hbar\omega(q)}2\leq e(q)
 \leq\frac{\hbar\omega(q)}2
       +\frac{9\Gamma\hbar^2}{32\mathcal V_3\omega(q)^2},
 \qquad e_1(q)\geq\frac{3\hbar\omega(q)}2.
\end{equation}
In the explicit parameter regime
\begin{equation}\label{fq:regime}
 \frac{9\Gamma\hbar}{32\mathcal V_3\omega_*^3}\leq\frac12,
 \quad\hbox{equivalently}\quad
 \mathcal V_3\geq\frac{9\Gamma\hbar}{16\omega_*^3},
\end{equation}
these imply, simultaneously for every $q$,
\begin{equation}\label{fq:separation}
 e(q)\leq\hbar\omega(q),\qquad
 e_1(q)-e(q)\geq\lambda(q):=\frac{\hbar\omega(q)}2>0.
\end{equation}
\end{lemma}
\begin{proof}
Compactness gives existence of a bottom eigenfunction. The scalar
kinetic form does not increase on taking its modulus, so a
nonnegative real eigenfunction can be selected. Alternatively, the
continuous nonnegative potential has the Brownian-bridge kernel
\[
 \frac{\exp[-(y-z)^2/(4a_y\tau)]}
 {\sqrt{\mathcal V_3/2}\sqrt{4\pi a_y\tau}}
 \mathbb E\exp\left[-\int_0^\tau
    (b_qZ_s^2+cZ_s^4)\,ds\right]>0\qquad(\tau>0).
\]
Here $Z$ is the bridge from $y$ to $z$ with diffusion covariance
$2a_y$. Each path is continuous with bounded image, so its
exponential is strictly positive and bounded by one. The formula
follows first for bounded truncations by the convergent Duhamel
series and then by increasing-form resolvent convergence and
dominated convergence of the bridge exponential. Consequently the
semigroup is positivity improving. The modulus of any bottom
eigenfunction is again a bottom eigenfunction; if a real bottom
eigenfunction changed sign, positivity improvement on both its
positive and negative parts would imply
$|e^{-\tau h(q)}u|<e^{-\tau h(q)}|u|$, contradicting its eigenvalue
equations. Orthogonalization against a positive bottom eigenfunction
therefore rules out a second independent real one. Reality gives
complex dimension one. The eigenfunction is smooth by the
one-dimensional eigenvalue equation. A zero of a nonnegative smooth
solution would have derivative zero, and uniqueness for that
second-order ODE would force the solution to vanish. Thus the
selected eigenfunction is strictly positive everywhere.

Put $\beta(q)=\mathcal V_3\omega(q)/(2\hbar)$ and use the exactly
normalized harmonic comparison vector
\[
 \chi_q(y)=\bigl(\mathcal V_3/2\bigr)^{-1/4}
       \bigl(\beta(q)/\pi\bigr)^{1/4}
       \exp[-\beta(q)y^2/2],\qquad
 \|\chi_q\|_{\mathscr K}=1.
\]
The harmonic part $h_{\mathrm{osc}}(q)=-a_y\partial_y^2+b_qy^2$
factorizes as
\[
 h_{\mathrm{osc}}(q)
 =a_y(-\partial_y+\beta y)(\partial_y+\beta y)
       +\hbar\omega/2.
\]
Its ground state is $\chi_q$ and its first excitation has energy
$3\hbar\omega/2$. To check the latter without presuming a missing
spectral level, the harmonic operator has compact resolvent; its
lowering operator $\partial_y+\beta y$ maps any eigenfunction above
the ground to an $L^2$ distributional eigenfunction with energy
lowered by $\hbar\omega$. Its norm is nonzero by the factorization.
No energy strictly between $\hbar\omega/2$ and $3\hbar\omega/2$
is possible. The function $y\chi_q$ attains the latter energy.
The operator-core fact above identifies the distributional
eigenfunctions with the self-adjoint operator domain.

Since $cy^4\geq0$, the variational principle and its two-dimensional
min--max version give the lower bounds for $e$ and $e_1$.
Gaussian integration gives, in the retained measure,
\[
 \langle\chi_q,y^2\chi_q\rangle=\frac{\hbar}{\mathcal V_3\omega},
 \qquad
 \langle\chi_q,y^4\chi_q\rangle
       =\frac{3\hbar^2}{\mathcal V_3^2\omega^2}.
\]
Its Rayleigh quotient is the upper bound in
\eqref{fq:energy-bounds}. This is only a trial-vector bound:
$h(q)\chi_q=(\hbar\omega/2+cy^4)\chi_q$ is not proportional
to $\chi_q$ because $c>0$. In particular the true state
$\Omega_q$ is not replaced by this Gaussian.
Finally $\omega(q)\geq\omega_*$ and \eqref{fq:regime} make the
upper correction at most $\hbar\omega(q)/2$, proving
\eqref{fq:separation}.
\end{proof}

\begin{lemma}[An exact smooth eigenline, with its phase fixed]
\label{fq:bundle}
The functions $q\mapsto e(q)$ and $q\mapsto\Omega_q\in\mathscr K$
are locally real analytic, the latter also in the common local
operator graph norm. Thus they define a global normalized positive
ground-state line over the entire original $q$-axis. In this phase,
\[
 \langle\Omega_q,\partial_q\Omega_q\rangle=0,
 \qquad \mathcal A(q):=i\langle\Omega_q,\partial_q\Omega_q\rangle=0.
\]
Put $P_q=|\Omega_q\rangle\langle\Omega_q|$, $Q_q=I-P_q$, and
\[
 R_q=(h(q)-e(q))^{-1}Q_q,\qquad
 d:=\frac{3\Gamma\mathcal V_3}{2}.
\]
The inverse is on the orthogonal complement of the actual ground
state. The exact derivative identities are
\begin{align}
 e'(q)&=dq\,\langle\Omega_q,y^2\Omega_q\rangle,\notag\\
 \Omega'_q&=-dq\,R_qQ_qy^2\Omega_q,\label{fq:first-derivative}\\
 Q_q\Omega''_q
 &=-R_qQ_q\bigl[d y^2\Omega_q
                     +2(dq y^2-e'(q))\Omega'_q\bigr],\notag\\
 \langle\Omega_q,\Omega''_q\rangle&=-\|\Omega'_q\|^2.\notag
\end{align}
These identities hold without imposing \eqref{fq:regime}.
\end{lemma}
\begin{proof}
For $u$ in the operator domain of a fixed fibre $h(q_0)$,
\[
 c\|y^2u\|^2\leq\langle u,h(q_0)u\rangle
                  \leq\|u\|\,\|h(q_0)u\|.
\]
The arithmetic--geometric mean inequality implies that multiplication
by $y^2$ is $h(q_0)$-bounded with arbitrarily small relative bound.
Since $h(q)-h(q_0)=(b_q-b_{q_0})y^2$, nearby real fibres have the
same operator domain. One can see the corresponding resolvent
analyticity directly: at a fixed resolvent point $z$,
\[
 (h(q)-z)^{-1}=(h(q_0)-z)^{-1}
 \bigl[I+(b_q-b_{q_0})y^2(h(q_0)-z)^{-1}\bigr]^{-1},
\]
where the last inverse is a locally convergent Neumann series.
For real $q$, the closed symmetric perturbation is self-adjoint:
the same formula gives surjectivity at two conjugate nonreal
points chosen far enough from the real axis. It agrees with the
Friedrichs realization on the compact smooth operator core.
The displayed formula is analytic also as a map into the original
operator domain equipped with its graph norm.

A small contour around the isolated simple eigenvalue $e(q_0)$
defines the spectral projection by the resolvent integral.
It is locally analytic in that graph norm, remains of rank one
because close projections have isomorphic ranges, and gives the
locally analytic eigenvalue. For a fully fixed positive reference
vector use $\chi_*(y)=\chi_0(y)$, the preceding Gaussian at
$\omega_*$. Then the actual state is exactly
\begin{equation}\label{fq:exact-state}
 \Omega_q=\frac{P_q\chi_*}{\|P_q\chi_*\|_{\mathscr K}}.
\end{equation}
The denominator is nonzero because the ground state and $\chi_*$
are both strictly positive. This formula selects the same real
positive vector on overlaps of the local contour constructions,
so the local analytic sections give a global section.
It can equivalently be obtained by the convergent normalized
heat iteration $e^{-\tau h(q)}\chi_*/\|e^{-\tau h(q)}\chi_*\|$
as $\tau\to\infty$, by the discrete spectral expansion and the
nonzero ground-state overlap.

The real normalization identity gives
$\langle\Omega_q,\Omega'_q\rangle=0$. Differentiating the
eigenvalue equation in its common graph domain gives
$(h-e)\Omega'=-(dq y^2-e')\Omega$. Taking its ground-state
component proves the first identity, and taking its complementary
component proves the second. A second differentiation gives
$(h-e)\Omega''=-(dy^2-e'')\Omega-2(dq y^2-e')\Omega'$.
Its complementary part gives the asserted reduced-resolvent
formula; differentiating the normalization twice gives its
ground-state component. All products in these identities are
in $\mathscr K$, since $y^2$ is bounded in the local graph norm.
\end{proof}

\subsection{Explicit global derivative bounds for the true quartic state}

\begin{lemma}[Moment and reduced-resolvent estimates]
\label{fq:derivative-bounds}
Assume the proved parameter regime \eqref{fq:regime}. Then, globally
in $q$,
\begin{align}
 \|\Omega'_q\|&\leq M_1:=\frac{\sqrt{15\Gamma}}{\omega_*},
       \label{fq:M1}\\
 \|Q_q\Omega''_q\|&\leq M_2:=
 \frac{18\sqrt5\,\Gamma}{\omega_*^2}
 +\frac{192}{5^{3/4}}
        \sqrt{\frac{\Gamma\mathcal V_3}{\hbar\omega_*}}.
       \label{fq:M2}
\end{align}
\end{lemma}
\begin{proof}
In this proof all fibre inner products use the retained measure
$d\mu_y$. Let $\mu_j=\langle\Omega_q,y^j\Omega_q\rangle$ for
even $j$. The energy identity gives $b_q\mu_2\leq e(q)$.
To get the necessary sharper fourth moment, take the bounded
Lipschitz multiplier $F_R(y)=\max(-R,\min(y,R))$.
Testing the weak eigenvalue equation with $F_R^2\Omega_q$ gives
the exact identity
\[
 a_y\|(F_R\Omega_q)'\|^2
 +\int(b_qy^2+cy^4)F_R^2|\Omega_q|^2\,d\mu_y
 =e(q)\|F_R\Omega_q\|^2
          +a_y\int|F_R'|^2|\Omega_q|^2\,d\mu_y.
\]
The test lies in the form domain because $F_R$ and its derivative
are bounded. Since $|F_R|\leq|y|$, $|F_R'|\leq1$, dropping
nonnegative terms and letting $R\to\infty$ by Fatou yields
\[
 b_q\mu_4\leq e(q)\mu_2+a_y,
 \qquad \mu_4\leq\frac{e(q)^2}{b_q^2}+\frac{a_y}{b_q}.
\]
By \eqref{fq:separation}, this gives the explicit estimates
\begin{equation}\label{fq:moments}
 \mu_2\leq\frac{4\hbar}{\mathcal V_3\omega(q)},\qquad
 \|y^2\Omega_q\|=\sqrt{\mu_4}
       \leq\frac{2\sqrt5\,\hbar}{\mathcal V_3\omega(q)}.
\end{equation}
No Gaussian moment has been assigned to the quartic ground state.

From the proved gap, $\|R_q\|\leq1/\lambda(q)$. If $v\in\mathscr K$
and $w=R_qQ_qv$, the exact equation $(h-e)w=Q_qv$ and $w\perp\Omega_q$
give
\[
 c\|y^2w\|^2\leq\langle w,h(q)w\rangle
 =e(q)\|w\|^2+\operatorname{Re}\langle w,v\rangle
 \leq\frac{e(q)+\lambda(q)}{\lambda(q)^2}\|v\|^2.
\]
Thus \eqref{fq:first-derivative} implies the pointwise estimates
\begin{align}
 \|\Omega'_q\|&\leq
       L(q):=\frac{6\sqrt5\,\Gamma|q|}{\omega(q)^2},\notag\\
 \|y^2\Omega'_q\|&\leq
       L(q)\sqrt{\frac{e(q)+\lambda(q)}c}
       \leq4L(q)\sqrt{\frac{\hbar\omega(q)}{\Gamma\mathcal V_3}},
       \label{fq:pointwise-first}\\
 |e'(q)|&\leq\frac{6\Gamma\hbar|q|}{\omega(q)}.\notag
\end{align}
Substitution in the exact second-derivative identity gives
\begin{equation}\label{fq:pointwise-second}
 \|Q_q\Omega''_q\|\leq
 \frac{6\sqrt5\,\Gamma}{\omega(q)^2}
 +144\sqrt5\,\Gamma^{3/2}\sqrt{\frac{\mathcal V_3}{\hbar}}
                      \frac{q^2}{\omega(q)^{5/2}}
 +144\sqrt5\,\Gamma^2\frac{q^2}{\omega(q)^4}.
\end{equation}
For completeness the exact scalar maxima used to make these
bounds global are
\begin{align*}
 \sup_q\frac{|q|}{\omega_*^2+3\Gamma q^2}
       &=\frac1{2\sqrt{3\Gamma}\,\omega_*},\\
 \sup_q\frac{q^2}{(\omega_*^2+3\Gamma q^2)^2}
       &=\frac1{12\Gamma\omega_*^2},\\
 \sup_q\frac{q^2}{(\omega_*^2+3\Gamma q^2)^{5/4}}
       &=\frac4{3\Gamma\,5^{5/4}\sqrt{\omega_*}}.
\end{align*}
Differentiation in $|q|$ or $q^2$ shows that the maxima occur,
respectively, at $q^2=\omega_*^2/(3\Gamma)$,
$q^2=\omega_*^2/(3\Gamma)$, and $q^2=4\omega_*^2/(3\Gamma)$;
each function decreases to zero at infinity. Applying these
maxima to \eqref{fq:pointwise-first} and
\eqref{fq:pointwise-second} gives \eqref{fq:M1} and
\eqref{fq:M2}, including the constants $18\sqrt5$ and
$192/5^{3/4}$.
\end{proof}

\subsection{The exact projected Hamiltonian and its off-diagonal part}

\begin{theorem}[Exact isometric correspondence, not an invariant product]
\label{fq:projection}
For all of the retained positive parameters, without the extra
restriction \eqref{fq:regime}, define the isometry
\[
 J:\mathscr L\longrightarrow\mathscr H_2,\qquad
 (Jf)(q,y)=f(q)\Omega_q(y),\qquad
 (J^*\Psi)(q)=\langle\Omega_q,\Psi(q,\cdot)\rangle_{\mathscr K}.
\]
Its range projection is the fibrewise projection
$(\mathsf P\Psi)(q)=P_q\Psi(q)$, and $\mathsf Q=I-\mathsf P$.
The exact pullback of the full quadratic form is the Friedrichs
form of
\begin{equation}\label{fq:effective}
 H_{\mathrm{eff}}=-\frac{\hbar^2}{2\mathcal V_3}\partial_q^2
       +V_q(q)+e(q)+B(q),\qquad
 B(q)=\frac{\hbar^2}{2\mathcal V_3}\|\Omega'_q\|_{\mathscr K}^2.
\end{equation}
Its domain is precisely
\[
 \mathcal D_{\mathrm{eff}}=
 \left\{f\in H^1(\mathbb R,d\mu_q):
  \int(V_q+e+B)|f|^2\,d\mu_q<\infty\right\},
\]
and $J^{-1}(\mathcal D(H_2^{1/2})\cap\operatorname{Ran}J)
=\mathcal D_{\mathrm{eff}}$.
The Berry connection is zero in the fixed positive real phase,
but $B(0)=0$ and $B(q)>0$ for every $q\ne0$.

For every $f\in C_c^\infty(\mathbb R)$, $Jf\in D(H_2)$ and
the exact operator identity is
\begin{align}
 H_2Jf&=JH_{\mathrm{eff}}f+\mathcal Rf,\label{fq:intertwining}\\
 (\mathcal Rf)(q)&=-a\bigl(2f'(q)\Omega'_q
                                  +f(q)Q_q\Omega''_q\bigr),
 \qquad \mathsf P\mathcal Rf=0.\notag
\end{align}
Thus $J^*H_2J=H_{\mathrm{eff}}$ on this core, while
$\mathsf QH_2J=\mathcal R$ is the retained off-diagonal residual.
The range of $J$ is not a reducing subspace.

Under \eqref{fq:regime}, these identities extend to all
$f\in D(H_{\mathrm{eff}})$, and
\begin{equation}\label{fq:residual}
 \|\mathcal Rf\|\leq
  2\sqrt a\,M_1\|H_{\mathrm{eff}}^{1/2}f\|
        +aM_2\|f\|,
 \qquad 0\leq B(q)\leq\frac{15\Gamma\hbar^2}
                                 {2\mathcal V_3\omega_*^2}.
\end{equation}
\end{theorem}
\begin{proof}
Normalization of the actual fibre state and Fubini give
$\|Jf\|=\|f\|$ and the displayed adjoint and projection.
For compact smooth $f$,
\[
 \partial_q(Jf)=f'\Omega_q+f\Omega'_q,\qquad
 \|\partial_q(Jf)(q)\|^2=|f'|^2+|f|^2\|\Omega'_q\|^2,
\]
because $\langle\Omega_q,\Omega'_q\rangle=0$. The fibre energy
is exactly $e(q)|f(q)|^2$, not a harmonic replacement. These
identities prove equality of the pulled-back forms on the core.
Core approximation gives the forward domain inclusion. Conversely,
if $Jf$ lies in the full form domain, local Hilbert-valued Sobolev
differentiation gives
$f'=\langle\Omega_q,\partial_q(Jf)\rangle$ distributionally;
the term with $\Omega'_q$ vanishes by normalization. Hence $f\in H^1$,
and the displayed nonnegative derivative and fibre-energy identities
give every required integral in $\mathcal D_{\mathrm{eff}}$.
This proves equality of the domains and closed forms. The positive
smooth effective potential defines its self-adjoint Friedrichs
operator with the operator core proved above. It has compact
resolvent because it is bounded below by $V_q(q)$, which tends
to infinity, and its kinetic coefficient $a$ is strictly positive.

By local graph regularity of $\Omega_q$, $Jf$ is in the full form
domain and its distributional Hamiltonian is in $\mathscr H_2$.
It therefore lies in $D(H_2)$. Expanding $\partial_q^2(f\Omega_q)$
and using $\langle\Omega_q,\Omega''_q\rangle=-\|\Omega'_q\|^2$
proves \eqref{fq:intertwining} and the orthogonality of its residual.

At $q=0$, \eqref{fq:first-derivative} gives $\Omega'_0=0$.
If $q\ne0$ and $\Omega'_q=0$, that identity and injectivity of
the reduced inverse on $Q_q\mathscr K$ would give
$Q_qy^2\Omega_q=0$. It would follow that $y^2\Omega_q$ is a
scalar multiple of $\Omega_q$. Multiplication by $y^2$ has no
nonzero $L^2$ eigenfunction, because every level set has measure
zero. Thus $\Omega'_q\ne0$ for $q\ne0$ and $B(q)>0$ there.
If the range of $J$ were reducing, its residual would vanish for
every compact smooth $f$. Since its coefficients are continuous
in $q$, choose at any $q_0\ne0$ a test function with
$f(q_0)=0$, $f'(q_0)=1$. The residual at $q_0$ is then
$-2a\Omega'_{q_0}\ne0$, a contradiction.

Under \eqref{fq:regime}, the global derivative estimates and the
triangle inequality give
\[
 \|\mathcal Rf\|\leq a(2M_1\|f'\|+M_2\|f\|).
\]
All terms in the effective potential are nonnegative, so
$\sqrt a\|f'\|\leq\|H_{\mathrm{eff}}^{1/2}f\|$.
This proves \eqref{fq:residual} and defines $\mathcal R$ as a
bounded map from the effective form domain with its form norm
to $\mathscr H_2$. For $f\in D(H_{\mathrm{eff}})$, choose an
operator-core sequence. It converges also in form norm by the
spectral inequality
$\|H_{\mathrm{eff}}^{1/2}u\|^2\leq\|u\|\|H_{\mathrm{eff}}u\|$.
The right side of \eqref{fq:intertwining} then converges in
$\mathscr H_2$. Closedness of $H_2$ shows that
$Jf\in D(H_2)$ and proves the identity on the full effective
operator domain. The bound on $B$ follows directly from $aM_1^2$.
\end{proof}

\subsection{Controlled physical-time and spectral comparisons}

\begin{theorem}[Explicit finite-energy error in a proved parameter regime]
\label{fq:dynamics}
Assume \eqref{fq:regime}. Let $f\in\mathcal D_{\mathrm{eff}}$
have $\|f\|=1$ and effective energy expectation
$\|H_{\mathrm{eff}}^{1/2}f\|^2\leq E$. For every physical time
$t\in\mathbb R$,
\begin{align}
 &\left\|e^{-itH_2/\hbar}Jf
                      -Je^{-itH_{\mathrm{eff}}/\hbar}f\right\|
       \leq |t|\,\mathcal E(E),\label{fq:time-error}\\
 &\mathcal E(E):=
 \frac{\sqrt{30\Gamma E}}{\omega_*\sqrt{\mathcal V_3}}
 +\frac{9\sqrt5\,\Gamma\hbar}{\mathcal V_3\omega_*^2}
 +\frac{96}{5^{3/4}}
             \sqrt{\frac{\Gamma\hbar}{\mathcal V_3\omega_*}}.
       \notag
\end{align}
The same bound holds for the leakage norm
$\|\mathsf Qe^{-itH_2/\hbar}Jf\|$.
In particular, for any prescribed $T>0$, $0<\rho<1$, the entirely
explicit sufficient volume regime
\begin{equation}\label{fq:target-volume}
 \mathcal V_3\geq\max\left\{
 \frac{9\Gamma\hbar}{16\omega_*^3},\
 \frac{270\Gamma E T^2}{\rho^2\omega_*^2},\
 \frac{27\sqrt5\,\Gamma\hbar T}{\rho\omega_*^2},\
 \frac{288^2\Gamma\hbar T^2}{5^{3/2}\rho^2\omega_*}
 \right\}
\end{equation}
makes both norms at most $\rho$ simultaneously for $|t|\leq T$.
All compared Hamiltonians still contain the original quartic
interaction, and no coupling has been set to one.
\end{theorem}
\begin{proof}
First take $f\in D(H_{\mathrm{eff}})$. The established operator-domain
map permits differentiation of
$e^{-i(t-s)H_2/\hbar}Je^{-isH_{\mathrm{eff}}/\hbar}f$.
Integrating its derivative gives the exact Duhamel identity
\[
 e^{-itH_2/\hbar}Jf-Je^{-itH_{\mathrm{eff}}/\hbar}f
 =-\frac{i}{\hbar}\int_0^t
    e^{-i(t-s)H_2/\hbar}\mathcal R
                       e^{-isH_{\mathrm{eff}}/\hbar}f\,ds.
\]
Unitarity and preservation of the effective form norm give the bound
$|t|(2\sqrt aM_1\sqrt E+aM_2)/\hbar$. Substitution of the exact
$a,M_1,M_2$ yields \eqref{fq:time-error}. For a general form-domain
vector, effective spectral cutoffs converge both in norm and form
norm and give the same estimate by continuity. The comparison
vector on the right lies in $\operatorname{Ran}J$, so applying
the contraction $\mathsf Q$ proves the leakage bound.
The last three lower bounds on volume in \eqref{fq:target-volume}
make respectively the three terms of $T\mathcal E(E)$ at most
$\rho/3$; the first is the separately proved fibre-gap regime.
Squaring the first and third nonnegative inequalities and directly
solving the second gives exactly their displayed constants.
\end{proof}

\begin{corollary}[An explicit approximate-eigenstate relation]
\label{fq:spectral-error}
In \eqref{fq:regime}, every normalized effective eigenvector
$H_{\mathrm{eff}}f=\Lambda f$ satisfies
\[
 \|(H_2-\Lambda)Jf\|\leq
       2\sqrt aM_1\sqrt\Lambda+aM_2
       =\hbar\mathcal E(\Lambda),\qquad
 \operatorname{dist}(\Lambda,\sigma(H_2))
       \leq\hbar\mathcal E(\Lambda).
\]
\end{corollary}
\begin{proof}
The effective operator has compact resolvent as proved above, so
its eigenvectors are in its operator domain, and $\Lambda>0$
because $e(q)\geq\hbar\omega_*/2$. The exact intertwining identity
and residual estimate prove the norm bound. The spectral theorem
for $H_2$ gives
$\|(H_2-\Lambda)u\|\geq\operatorname{dist}(\Lambda,\sigma(H_2))\|u\|$
for every operator-domain vector. Apply it to the normalized $Jf$.
\end{proof}

The preceding comparison is an actual isometry, an exact compressed
Hamiltonian, and an explicit finite-time error estimate for this
specific interacting two-coordinate mechanics. At fixed
$\hbar,\Gamma,L_1,m^2,E,T$, its stated error tends to zero as
$\mathcal V_3\to\infty$, for any family of initial states satisfying
the stated energy bound. This is not a volume-uniform mass-gap
claim and not a continuum limit. The estimate is conservative;
in particular the second-derivative bound deliberately uses the
positive quartic energy to control $y^2R_q$, which is sufficient
for the proved $\mathcal V_3^{-1/2}$ error without an unproved
adiabatic expansion or omission of $y^4$.

\paragraph{Context and attribution.}
Moving spectral subspaces, the Berry connection, diagonal
Born--Huang energy, and errors in time-dependent Born--Oppenheimer
comparisons are established concepts; see G.~Panati, H.~Spohn,
and S.~Teufel, \emph{The time-dependent Born--Oppenheimer approximation},
ESAIM: M2AN \textbf{41} (2007), 297--314,
\url{https://arxiv.org/abs/0712.4369},
\url{https://doi.org/10.1051/m2an:2007023}.
No new general approximation theorem is claimed here. The spectral
bundle, every retained coefficient, and the conservative residual
and time estimates used above were proved directly for
\eqref{fq:H}. The bounded-below bridge formula used in the positivity
argument is also stated in B.~Simon, \emph{A Feynman--Kac Formula for
Unbounded Semigroups}, Theorem~1.1,
\url{https://arxiv.org/abs/math-ph/9907022}; its general unbounded-below
extension is not needed here.
